\chapter{The Conservation of Topology ($N$) and Time Dilation}
\label{app:time_dilation}
\addtocontents{toc}{\protect\setcounter{tocdepth}{0}}

\noindent\textbf{Abstract:} Standard cosmology utilizes Supernova Time Dilation as definitive proof of expanding space. Tension Field Theory (TFT) mathematically guarantees this exact time dilation without metric expansion by enforcing the topological conservation of the wave packet.

\section{The Topological Invariant ($N$)}
A photon is a finite soliton wave characterized by a strict integer count of phase cycles ($N$). While the kinetic energy of the photon is not conserved (as it bleeds heat into the lattice), this topological integrity cannot be destroyed by linear propagation. The total cycle count is the product of the wave's frequency ($f$) and its temporal duration ($\Delta t$):
\begin{equation}
    N = f_{\mathrm{emit}} \cdot \Delta t_{\mathrm{emit}}
\end{equation}

\section{The Frequency Drop}
As derived in \hyperref[app:vacuumdamping_cosmicredshift]{Appendix \ref*{app:vacuumdamping_cosmicredshift}}, visco-elastic damping causes the wave's energy to decay exponentially (depositing the deficit as CMB phonons), which manifests directly as a drop in the observed frequency ($f_{\mathrm{obs}} < f_{\mathrm{emit}}$), related by the redshift factor:
\begin{equation}
    f_{\mathrm{obs}} = \frac{f_{\mathrm{emit}}}{1+z}
\end{equation}

\section{The Derivation of Time Dilation}
Because the topological cycle count must be strictly conserved upon observation ($N_{\mathrm{obs}} = N_{\mathrm{emit}}$), we equate the emitted and observed states:
\begin{equation}
    f_{\mathrm{obs}} \cdot \Delta t_{\mathrm{obs}} = f_{\mathrm{emit}} \cdot \Delta t_{\mathrm{emit}}
\end{equation}

Substitute the redshifted frequency ($f_{\mathrm{obs}}$):
\begin{equation}
    \left( \frac{f_{\mathrm{emit}}}{1+z} \right) \cdot \Delta t_{\mathrm{obs}} = f_{\mathrm{emit}} \cdot \Delta t_{\mathrm{emit}}
\end{equation}

Cancel the emission frequency ($f_{\mathrm{emit}}$) and solve for the observed duration ($\Delta t_{\mathrm{obs}}$):
\begin{equation}
    \Delta t_{\mathrm{obs}} = \Delta t_{\mathrm{emit}} (1+z)
\end{equation}

\vspace{0.5em}
\noindent\textbf{Conclusion:} The physical duration of the wave packet must proportionally increase to accommodate the lower frequency while conserving the cycle count ($N$). A distant supernova appears to happen in slow motion ($\Delta t_{\mathrm{obs}} > \Delta t_{\mathrm{emit}}$) strictly because the wave packet has physically elongated to conserve its topology against the cosmic damping. No expanding metric is required.
